\chapter{Сравнение методов}

Сравнение проводится по двум группам метрик: вычислительной эффективности (FLOPs, число параметров, время инференса) и качеству синтеза (PSNR, SSIM, LPIPS).

\section{Экспериментальная установка}

\subsection*{Модели для сравнения}

В рамках исследования были обучены и протестированы следующие модели на основе архитектуры IFRNet:
\begin{itemize}
	\item \texttt{IFRNet (baseline)}~---~оригинальная архитектура, предложенная в работе \cite{IFRNet2022}. 
	Использует полное цветовое пространство RGB и стандартную функцию потерь: комбинацию L1, Ternary и дистилляционного лосса.
	
	\item \texttt{IFRNet\_RGB2}~---~модификация с равномерным сокращением числа каналов в два раза на всех уровнях пирамиды признаков (каналы 16, 24, 36, 48 вместо 32, 48, 72, 96). 
	Функция потерь идентична базовой модели.
	
	\item \texttt{IFRNet\_YUV420}~---~модификация с переходом к цветовому пространству YUV420 и интерполяцией на основе яркостной компоненты $Y$. 
	Цветоразностные компоненты $U$ и $V$ переносятся с использованием оптического потока, вычисленного по яркости. 
	Функция потерь аналогична базовой модели.
	
	\item \texttt{IFRNet\_RGB2+LPIPS}~---~архитектура \texttt{IFRNet\_RGB2} с добавлением перцептуальной функции потерь LPIPS~\cite{zhang2018} в общую функцию потерь.
	
	\item \texttt{IFRNet\_YUV420+LPIPS}~---~архитектура \texttt{IFRNet\_YUV420} с добавлением перцептуальной функции потерь LPIPS.
\end{itemize}

\subsection*{Метрики оценки}

Оценка производилась по двум категориям метрик:

\textbf{Метрики эффективности:}
\begin{itemize}
	\item \textbf{FLOPs} (Floating Point Operations)~---~количество операций с плавающей запятой, необходимое для обработки одного кадра. Рассчитывалось с помощью библиотеки \texttt{thop} для изображений размером $1920 \times 1088$.
	\item \textbf{Params}~---~количество обучаемых параметров модели (в миллионах).
	\item \textbf{Inference Time}~---~среднее время инференса на один кадр (в секундах).
\end{itemize}

\textbf{Метрики качества:}
\begin{itemize}
	\item \textbf{PSNR} (Peak Signal-to-Noise Ratio)~---~пиксельная метрика, измеряющая отношение сигнал/шум в децибелах.
	\item \textbf{SSIM} (Structural Similarity Index)~---~метрика структурного сходства, учитывающая яркость, контраст и структуру изображения~\cite{wang2004}.
	\item \textbf{LPIPS} (Learned Perceptual Image Patch Similarity)~---~перцептуальная метрика, использующая признаки из предобученной сети VGG для оценки визуального сходства~\cite{zhang2018}.
\end{itemize}

\subsection*{Датасет и условия обучения}

Обучение всех моделей проводилось на датасете Vimeo90K. Входные изображения имели разрешение $1920 \times 1088$ пикселей. 
Для моделей с LPIPS в функции потерь использовался коэффициент $\lambda = 0.1$ для взвешивания перцептуальной компоненты.

\section{Результаты экспериментов}

\subsection*{Количественное сравнение}

Результаты сравнения всех моделей представлены в таблице~\ref{tab:comparison}.

\begin{table}[htbp]
	\centering
	\caption{Сравнение методов оптимизации IFRNet}
	\label{tab:comparison}
	\adjustbox{max width=\textwidth}{
		\begin{tabular}{l c c c c c c c}
			\toprule
			\textbf{Модель} & \texttt{FLOPs (↓)} & \texttt{Params (↓)} & \texttt{Inference Time (↓)} & \texttt{PSNR (↑)} & \texttt{SSIM (↑)} & \texttt{LPIPS (↓)} & \texttt{DISTS (↓)} \\
			& & & \textbf{(с/кадр)} & & & & \\
			\midrule
			\texttt{IFRNet (baseline)} & 479.441G & 4.960M & 0.034 & 35.698 & 0.9790 & 0.0517 & 0.0421 \\
			\texttt{IFRNet\_RGB2} & 126.191G & 1.00M & 0.015 & 34.897 & 0.9756 & 0.0542 & 0.0453 \\
			\texttt{IFRNet\_YUV420} & 122.682G & 1.00M & 0.017 & 34.446 & 0.9737 & 0.0622 & 0.0527 \\
			\texttt{IFRNet\_RGB2+LPIPS} & 126.191G & 1.00M & 0.015 & 34.659 & 0.9743 & 0.0380 & 0.0389 \\
			\texttt{IFRNet\_YUV420+LPIPS} & 122.682G & 1.00M & 0.017 & 35.121 & 0.9761 & 0.0481 & 0.0405 \\
			\bottomrule
		\end{tabular}
	}
\end{table}

\subsection*{Анализ эффективности}

\textbf{Сокращение вычислительной сложности.}
Все предложенные методы оптимизации демонстрируют значительное снижение вычислительной сложности по сравнению с базовой моделью IFRNet. 
Наиболее существенное ускорение достигается за счёт сокращения ширины сети: модели \texttt{IFRNet\_RGB2} и \texttt{IFRNet\_RGB2+LPIPS} показывают снижение FLOPs в 3.8 раза (с 479.441G до 126.191G) и ускорение инференса более чем в 2 раза (с 0.034 до 0.015 секунды на кадр).

Модель \texttt{IFRNet\_YUV420} также показывает значительное снижение FLOPs до 122.682G, что сопоставимо с RGB2-версиями, однако время инференса несколько выше (0.017 с/кадр), что может быть связано с дополнительными операциями преобразования цветовых пространств и деформации цветовых каналов.

\textbf{Число параметров.} 
Сокращение ширины сети в два раза позволило уменьшить число параметров с 4.96M до 1.00M. 
Модели на основе YUV420 сохраняют такое же количество параметров, что объясняется неизменной архитектурой сети (изменения касаются только входных данных и стратегии синтеза цветовых компонент).

\subsection*{Анализ качества}

\textbf{Пиксельные метрики.} 
Все оптимизированные модели показывают некоторое снижение PSNR и SSIM по сравнению с базовой архитектурой. 
Наименьшую потерю качества демонстрирует модель \texttt{IFRNet\_YUV420+LPIPS}, которая достигает PSNR = 35.121 дБ при снижении FLOPs в 3.9 раза. Это подтверждает гипотезу о том, что цветовая информация менее критична для восприятия, и её обработка с пониженным разрешением позволяет сохранить высокое качество синтеза.

Модели без LPIPS в функции потерь (\texttt{IFRNet\_RGB2} и \texttt{IFRNet\_YUV420}) показывают ожидаемое снижение PSNR на 0.8-1.25 дБ, что является приемлемой ценой за значительное ускорение работы модели.

\textbf{Перцептуальное качество.}
Наиболее интересные результаты получены по метрике LPIPS. 
Добавление перцептуальной функции потерь в процесс обучения позволило существенно улучшить визуальное качество генерируемых кадров. 
Модель \texttt{IFRNet\_RGB2+LPIPS} показывает значение LPIPS = 0.0380, что даже лучше, чем у базовой модели (0.0517), при том же уровне вычислительной эффективности, что и \texttt{IFRNet\_RGB2}.

Модель IFRNet\_YUV420+LPIPS также демонстрирует значительное улучшение LPIPS по сравнению с версией без LPIPS (0.0481 против 0.0622), при этом оставаясь близкой к базовой модели по этому показателю.

\subsection*{Сравнительный анализ баланса эффективности и качества}

На рисунке~\ref{fig:pareto} представлена визуализация компромисса между вычислительной эффективностью и качеством синтеза для всех рассмотренных моделей.

% \begin{figure}[htbp]
	%     \centering
	%     \includegraphics[width=\textwidth]{pareto_chart.png}
	%     \caption{График компромисса между FLOPs и качеством (LPIPS) для различных моделей}
	%     \label{fig:pareto}
	% \end{figure}

Можно выделить три группы моделей по соотношению качество/эффективность:

\begin{enumerate}
	\item \textbf{максимальная эффективность}: IFRNet\_RGB2 и IFRNet\_RGB2+LPIPS обеспечивают наибольшее ускорение (в 3.8 раза по FLOPs и в 2.3 раза по времени инференса). При этом модель с LPIPS значительно превосходит версию без LPIPS по перцептуальному качеству.
	
	\item \textbf{сбалансированное решение}: IFRNet\_YUV420+LPIPS предлагает оптимальный баланс между скоростью и качеством, обеспечивая значительное ускорение (FLOPs ↓ в 3.9 раз) при минимальной потере качества по пиксельным метрикам и улучшении перцептуальных характеристик.
	
	\item \textbf{базовая модель}: IFRNet (baseline) демонстрирует наилучшее пиксельное качество, но ценой значительно более высоких вычислительных затрат.
\end{enumerate}

\section{Визуальное сравнение}

Для качественной оценки результатов было проведено визуальное сравнение интерполированных кадров, полученных различными моделями. На рисунках~\ref{fig:visual_comparison1} и~\ref{fig:visual_comparison2} представлены примеры синтезированных кадров для различных сцен.

% \begin{figure}[htbp]
	%     \centering
	%     \includegraphics[width=\textwidth]{visual_comparison1.png}
	%     \caption{Визуальное сравнение интерполяции для сцены с быстрым движением}
	%     \label{fig:visual_comparison1}
	% \end{figure}

% \begin{figure}[htbp]
	%     \centering
	%     \includegraphics[width=\textwidth]{visual_comparison2.png}
	%     \caption{Визуальное сравнение интерполяции для сцены со сложными текстурами}
	%     \label{fig:visual_comparison2}
	% \end{figure}

\textbf{Анализ визуальных результатов:}
\begin{itemize}
	\item Модели без LPIPS (IFRNet\_RGB2, IFRNet\_YUV420) демонстрируют некоторое сглаживание текстур и потерю мелких деталей, что особенно заметно на сложных природных сценах.
	\item Добавление LPIPS в функцию потерь позволяет восстановить резкость и текстуры: модели с LPIPS дают более детализированные и естественные изображения.
	\item Модель IFRNet\_YUV420+LPIPS показывает результаты, визуально близкие к базовой модели, при значительно меньшей вычислительной нагрузке.
\end{itemize}